% !TeX root = ../抽象代数笔记.tex
% 由主文档通过 \input 载入；请编译主文档。
\setcounter{chapter}{6} % 保留原书章号。
\chapter{西罗定理与有限群}\label{ch:sylow}

\begin{description}
  \item[学习目标] 利用素数幂阶子群的存在、共轭和数量限制研究有限群。
  \item[先修知识] 第\ref{ch:actions}章，尤其是轨道计数与共轭作用。
\end{description}

\section{柯西定理与素数幂子群}
证明柯西定理（Cauchy's theorem）：若素数 $p$ 整除有限群 $G$ 的阶，则 $G$ 含有一个 $p$ 阶元素。用循环移位作用给出一种证明，明确作用集合及其不动点。比较它与拉格朗日定理，说明其结论只处理素数阶而非任意因子阶。

\section{西罗定理的三个结论}
设 $|G|=p^a m$，其中 $a\geq1$ 且 $p\nmid m$。定义西罗 $p$ 子群（Sylow $p$-subgroup）为 $p^a$ 阶子群。完整陈述并证明：西罗 $p$ 子群存在；每个 $p$ 子群包含于某个西罗 $p$ 子群且所有西罗 $p$ 子群互为共轭；其数目 $n_p$ 满足 $n_p\mid m$ 与 $n_p\equiv1\pmod p$。

\section{正规化子与证明的组织}
对 $P\leq G$ 定义正规化子 $N_G(P)=\{g\in G:gPg^{-1}=P\}$。在共轭作用下推出 $n_p=[G:N_G(P)]$，并证明西罗 $p$ 子群正规当且仅当 $n_p=1$。将存在性、包含性、共轭性和数量公式分步证明，标明每一步使用的作用与计数条件。

\section{小阶群的分析与方法边界}
分析 $6$、$10$、$15$ 阶群的西罗子群数量，以 $15$ 阶群必循环说明数量限制、正规性与直积思想的结合。对 $12$ 阶群展示仅凭同余与整除条件可能保留多个候选值；进一步结论需要元素计数、置换作用或结构分析，不能把候选值当成已实现的群。

\section{例题、练习与自检}
安排判断给定阶数的群是否必有非平凡真正规子群、计算 $S_4$ 的西罗子群以及排除某些单群阶数的题目。自检要求能够说明“存在某个正规子群”与“完整分类全部群”之间的区别，并检查每次使用 $n_p\mid m$ 时 $m$ 的取值。

